\documentclass[11pt]{article}
\usepackage{amsmath,amssymb,amsthm,hyperref}
\newtheorem{theorem}{Theorem}
\newtheorem{lemma}[theorem]{Lemma}
\newtheorem{corollary}[theorem]{Corollary}
\newcommand{\Q}{\mathbb Q}
\newcommand{\C}{\mathbb C}
\newcommand{\CH}{\operatorname{CH}}
\newcommand{\cl}{\operatorname{cl}^{\mathrm{can}}}
\title{GST curvature: exact missing-mode sequence and completion of divisorial returns}
\author{Mathematics before Lean; author self-review}
\date{3 November 2026}
\begin{document}\maketitle

\section{Claim and checkpoint}
The starting revision is \texttt{e40e328b80d3ede78202234054891b1e56edf0b3}.
The original input is a connected smooth projective complex $n$-fold $X$,
$1\leq p\leq n$, and a prescribed rational Hodge class
$\alpha\in H^{2p}(X,\Q)\cap H^{p,p}(X)$.
The universal target is an actual $Z_\alpha\in\CH^p(X)_\Q$
with $\cl_X Z_\alpha=\alpha$. Components are handled separately.
The central-fold curvature theorem is already proved at hand level in its
stated family. Here we make its missing-mode exact sequence explicit in
the custom module, then solve the previously untreated \emph{divisorial}
boundary extraction, including its rational Hodge-lift issue.
Existence of a nullifying algebraic open for every $\alpha$ is not assumed.

\section{The curvature conclusion, with its full linear exact sequence}
For the actual central-fold family of the earlier derivation let
$L:\C^r\to\C^m$ be the transpose of the total integral curvature matrix.
Discarding the already pure base contribution, closed vertical degree-one
classes form
\[
 E_L=\{(u,v)\in\C^r\oplus\C^r:Lu=0\}.
\]
Define $q:E_L\to\operatorname{im}L$ by $q(u,v)=Lv$ and
\[
 P_L=\{(a+b,i(a-b)):a,b\in\ker L\}.
\]
Then the actual linear sequence is
\[
 0\longrightarrow P_L\longrightarrow E_L
   \xrightarrow{q}\operatorname{im}L\longrightarrow0.       \tag{1}
\]
Indeed $q$ is onto: use $(0,v)$ for any $Lv$.
Its kernel has $Lu=Lv=0$, with the explicit unique decomposition
$a=(u-iv)/2$, $b=(u+iv)/2$. Conversely that decomposition is in the
kernel. In particular $\dim(E_L/P_L)=\operatorname{rank}L$.
The geometric comparison and filtered representatives proved in
\texttt{GST\_whole\_space\_central\_causal\_fold\_2026\_11\_03.tex}
identify (1) with the actual degree-one opposedness quotient.
Consequently the stated curvature theorem is
\[
 \delta^1=\operatorname{rank}\Bigl(\sum_e K_e\Bigr).
                                                               \tag{2}
\]
In this particular lattice-defined family, zero total matrix gives the
actual projective product, so purity returns in \emph{all} degrees.
For arbitrary holomorphic torus torsors zero rational curvature alone
does not imply projectivity: their torsion-gluing criterion is retained.
Neither assertion is an original-$X$ universal Chow source theorem.

\section{The precise missing divisorial construction}
Let $f:Y\to X$ be a surjection with $Y$ smooth connected projective,
$\dim Y=n+r$. Fix a very ample $h\in\CH^1(Y)$ and the positive generic
degree $d$ with $f_*(h^r)=d[X]$. The previously constructed genuine return is
\[
 R_{\CH}(z)=d^{-1}f_*(h^r z),\qquad R_B f^*=1.             \tag{3}
\]
Let $D=\bigcup_i D_i$ be a reduced simple-normal-crossings divisor on $Y$,
$j:U=Y\setminus D\hookrightarrow Y$, and $i_i:D_i\hookrightarrow Y$.
Every $D_i$ is smooth projective; all intersections are transverse.

\begin{theorem}[Exact divisorial kernel, every degree]\label{thm:kernel}
For every $p\geq1$,
\[
 \ker j^*|_{H^{2p}(Y,\Q)}
 =\operatorname{im}\left(G:\bigoplus_i H^{2p-2}(D_i,\Q)(-1)
                \longrightarrow H^{2p}(Y,\Q)\right),\quad
 G((b_i))=\sum_i(i_i)_*b_i.                              \tag{4}
\]
If $a$ in this kernel is rational of type $(p,p)$, it has a rational
preimage with every $b_i$ of type $(p-1,p-1)$.
\end{theorem}
\begin{proof}
Write $D^{(s)}$ for the disjoint union of the $s$-fold component
intersections, with $D^{(0)}=Y$. The logarithmic weight spectral sequence,
in the indexing used here, is
\[
 E_1^{-s,t}=H^{t-2s}(D^{(s)},\Q)(-s)
       \Longrightarrow H^{t-s}(U,\Q).                    \tag{5}
\]
The first differential is the alternating sum of Gysin maps; in the
column adjacent to $Y$ it is exactly $G$. For fixed $t$ the terms are
pure of weight $t$ and the first differential preserves their Hodge
structures. Its cohomology is again pure of weight $t$.
A higher differential $d_a$ for $a\geq2$ changes $t$ to $t-a+1$,
so is zero as a morphism between unequal pure weights. Thus (5)
degenerates at $E_2$. There are no positive columns. The edge map from
$H^{2p}(Y)$ to $H^{2p}(U)$ therefore has kernel exactly the incoming
$d_1$ image in (4). Higher intersections impose relations on $G$;
they introduce no further incoming differential to this edge.

For completeness the rational Hodge lift requires more than selecting
the complex $(p-1,p-1)$ component of an arbitrary rational lift.
Set $B=\bigoplus_i H^{2p-2}(D_i,\Q)(-1)$ and $K=\ker G$.
Choose a rational polarization on the pure Hodge structure $B$.
Its restriction to $K$ is nondegenerate: positivity after the Weil
operator on the Hodge-stable real subspace rules out a radical.
Thus $B=K\oplus K^\perp$ is a \emph{rational Hodge} decomposition.
The restriction $G:K^\perp\to\operatorname{im}G$ is a rational Hodge
isomorphism. Apply its inverse to $a$. After undoing the Tate twist
the components have the asserted rational types. This proves the lift
without assuming any higher Hodge conjecture.
\end{proof}
The exact external inputs for (5) are Deligne, \emph{Th\'eorie de Hodge II}
(1971), \S3.1 residue description and \S3.2, especially (3.2.4),
(3.2.8)--(3.2.10) and (3.2.13)(iii). Rational polarizability is stated
in (2.2.6). The original published scan was checked at
\href{https://publications.ias.edu/sites/default/files/Number15.pdf}{IAS}.
This is the $E_1$ weight convention, not an incorrect assertion that
the differently indexed Leray sequence degenerates at $E_2$.

\begin{theorem}[Completed codimension-two boundary source]\label{thm:cycle}
For $p=2$, if the actual constructed open above satisfies
$j^*f^*\alpha=0$, actual algebraic line bundles $L_i$ on $D_i$ and a
common positive integer $N$ give the original-$X$ cycle
\[
 \boxed{Z_\alpha=\frac1{dN}f_*\left(h^r\cdot
                  \sum_i(i_i)_*c_1^{\CH}(L_i)\right),
       \qquad\cl_X Z_\alpha=\alpha.}                    \tag{6}
\]
Here the $L_i$ are constructed from $\alpha$ by the rational Hodge
lift in Theorem~\ref{thm:kernel}, denominator clearing, and the actual
exponential-sequence/GAGA construction, not assumed as cycle witnesses.
\end{theorem}
\begin{proof}
Take $a=f^*\alpha$. Formula (4) and its rational lift give rational
$(1,1)$ classes $b_i$ with $G(b)=a$. Choose one $N>0$ making all
$Nb_i$ integral modulo torsion. Each integral Hodge class maps to zero
in $H^2(D_i,\mathcal O_{D_i})$, so the analytic exponential exact
sequence constructs a holomorphic line $L_i$ with Chern class $Nb_i$.
Projective GAGA makes it algebraic. A nonzero rational section provides
an actual Cartier divisor if a cycle representative is desired.
Thus the sum of their proper pushforwards has canonical class $Na$.
Apply (3) and divide by $N$. Products, proper pushforward and Chern
classes commute with the canonical Betti comparison, proving (6).
There is no extra intersection correction to the sum of Gysin classes.
\end{proof}
This includes an arbitrary singular divisor after embedded resolution:
resolve the pair to a smooth projective $Y'$ with reduced SNC boundary,
apply (6) to $Y'\to X$, and push its actual cycle back. The open is
unchanged over its smooth complement. More generally a dominant map
from a smooth quasi-projective nullifying $U$ can first be compactified
by closure of its graph over projective $X$, then resolved relative to $U$.
It must be an \emph{actual algebraic} $U$ with the specified pullback;
the higher rational homotopy fiber does not supply it automatically.

For general $p$ the same proof completes the exact Gysin Hodge lift.
Actual cycles follow if the constructed $b_i$ have codimension-$(p-1)$
Chow representatives on the actual $D_i$. Only the $p=2$ case removes
that lower-degree source requirement by the unconditional $(1,1)$ theorem.

\section{Execute a direct arbitrary-input open construction}
Start again from the actual arbitrary $(X,p,\alpha)$ and choose a smooth
divisor $D\in|H^{\otimes m}|$ for $H$ very ample and $m>0$.
This constructs an actual affine open $U=X\setminus D$ without
assuming $\alpha$ algebraic. Put $h=c_1(H)$ and $L(b)=h\cup b$.
Its exact nullification test is
\[
 \boxed{\ker(H^{2p}(X,\Q)\to H^{2p}(U,\Q))
                   =L H^{2p-2}(X,\Q).}                 \tag{7}
\]
For $2p\leq n$, weak Lefschetz gives
$H^{2p-2}(X)\simeq H^{2p-2}(D)$; the Gysin projection formula
$i_*i^*b=mh\cup b$ proves (7). The relevant weak Lefschetz range is
$2p-2<n-1$. For $2p>n$, hard Lefschetz makes $L$ surjective, and
every $L$-multiple restricts to zero because $H^{\otimes m}|_U$ is
trivial. Hence both sides in (7) are all of $H^{2p}(X)$.
These uses have their usual rational, smooth-projective hypotheses.

For $2p\leq n$, the rational Hodge Lefschetz decomposition is
\[
 \alpha=h\beta+\alpha_{\rm prim},\qquad
 \alpha_{\rm prim}\in P^{2p}(X)=\ker L^{n-2p+1}.        \tag{8}
\]
The summands are rational Hodge summands and have zero intersection.
Thus this actual open kills the input exactly when
$\alpha_{\rm prim}=0$. In particular a nonzero primitive residual
survives \emph{every} smooth ample-divisor choice, regardless of $m$.
It cannot be eliminated by changing a connection or a de Rham primitive:
its restriction cohomology class is nonzero.

For $p=2$, the $\beta$ in (8) is rational $(1,1)$ and constructs an
actual rational divisor $E$ by exponential sequence and GAGA.
The explicit source produced by this attempt is
$c_1^{\CH}(H)\cdot E$, with class $h\beta$; the residual is exactly
$\alpha_{\rm prim}$. When $n\leq3$ the standard hard-Lefschetz/top-class
argument produces all codimension-two inputs, as already recorded in
October; that known result is not a fresh universal breakthrough.
For $n\geq4$ this attempt has now computed its precise success sector
and the surviving prescribed-input class.

\section{An actual singular-boundary control for the primitive channel}
The arbitrary-$X$ implication of this control is precise: Theorem~\ref{thm:cycle}
handles resolved divisorial geometry beyond the smooth-ample kernel (7).
The following tests that distinction with an actual source; it is not
evidence of universal source existence or a replacement of the input variety.

Let $X=Q^4\subset\mathbf P^5$ be a smooth quadric. Write $A,B$ for
the two maximal-plane classes, so $h^2=A+B$ and
$\gamma=A-B$ is nonzero primitive. Choose a point $v\in Q^4$ and
its tangent hyperplane section $D$. This is the cone with vertex $v$
over $Q^2\simeq\mathbf P^1\times\mathbf P^1$.
Its ruling planes of both families contain $v$. Blow up $v$ in $X$:
$f:Y\to X$. The reduced total boundary is $\widetilde D\cup E$,
where $E\simeq\mathbf P^3$, $\widetilde D$ is the projective-line
bundle resolving this cone, and their transverse intersection is the
quadric $Q^2\subset E$. Thus this is an actual SNC compactification
of $X\setminus D$.

Let $\pi:\widetilde D\to Q^2$ be the ruling projection and
$M=\pi^*\mathcal O_{Q^2}(1,-1)$. The inverse image of a line in
either ruling is the strict transform of the corresponding plane
through $v$, with multiplicity one. Consequently the actual Chow identity
is
\[
 \boxed{f_*\widetilde i_*c_1^{\CH}(M)=A-B.}             \tag{9}
\]
Both sides of (9) are cycles, not merely numerical coordinates.
Here $r=0,d=N=1$. The correction class $c_1(M)$ is an actual line
on the resolved boundary, chosen by the prescribed difference of
the two rulings. In contrast every smooth ample hyperplane complement
retains $A-B$ by (7). The singular boundary succeeds because it has
the needed additional divisor geometry. This rules out extending the
smooth-ample obstruction to all divisorial opens. These quadric cycles
are classical; the result is a constructive check of the general return.

\section{Adversarial review and continuation}
The curvature theorem (2) is established at its stated family scope.
The exact sequence (1), the SNC kernel and rational Hodge lift (4),
and actual divisor-line return (6) are hand proofs under their explicit
geometric inputs. The source-dependent step (5) has been checked against
the primary text. Mathematical review here is author self-review.
The Lean changes prove algebraic map consequences and the explicit
missing-mode surjection/kernel; compilation remains with the owner.

Three failure controls are essential. Projectivity is not inferred from
zero curvature for arbitrary torsors. A complex type component of an
arbitrary rational lift is not silently declared rational. A generic
smooth ample boundary is not asserted to nullify a prescribed primitive
input. Equation (7) is the executed construction test, not a new
algebraicity assumption.

The divisorial extraction subconstruction is completed in codimension two
for actual nullifying algebraic opens, and in all degrees up to the
explicit lower-degree source input. Universal construction of such an open
for the surviving primitive input remains unresolved. The smooth ample
recipe is closed for that residual by (7); singular selected/excess opens
and other genuinely different source geometries retain their open status.
\end{document}
